\section{账本的接力：继位、财政与地方行政}
\label{sec:high-qing-fiscal-order-succession}

雍正初年的朝廷首先面对的不是一个可以用「整顿」一词收束的问题，而是一连串须由不同机构、不同地方分别结算的关系。皇帝在北京继承了名义上的统治权，也继承了州县亏空、旗人供养、河工与漕运的开支，以及西北军情所催生的粮饷需要。银两若只停在户部的数字上，不能养活边军；一项征收若只停在谕旨里，也不能说明纳户在何日、以何种折色、向谁缴纳。因而，雍正即位之后的用人、密折和财政措施，应同看作把消息、责任与资源重新接到中枢的努力，而不是一部已经顺畅运转的机器。

\eventref{EQ-1726-HAOXIAN-GUIGONG}{耗羡归公与养廉银的安排（1726）}正落在这条接力线上。州县正额俸给不足，附加征收又长期嵌在地方办事与官员生活之中；改革试图将其中一部分收入纳入较可核查的省级收支，并以养廉银为官员提供制度化的用度。它改变的首先是责任的书写方式：督抚要汇总、分拨和报解，州县要把原本在地方关系中处理的征收名目交到更长的文书链中，户部因而有了追问亏空和留支的依据。

这并不意味着同一套账在每一省、每一县落得一样。谕旨和《大清会典事例》能够说明朝廷如何界定名目；户部题本、奏销册及省县钱粮档案才能分辨定额、报解、实收、留存和欠额。即使账册保存下来，银、钱、实物和折色也未必使用同一单位或价格背景。将耗羡归公直接写成普遍减负，或写成地方附加从此消失，都会把制度的目标误当成民间已经经历的结果。对纳户、书吏、乡绅和地方官来说，更紧的核算可能减少某些任意性，也可能把催征、垫赔和解释责任转移到新的环节；哪些结果在何处发生，必须由具体地方的账簿、诉讼和契约追问。

\subsection{继承的是机构，也继承未清的账}

雍正的继位不能从后来围绕传位的故事倒推其最初数日的全部内情。实录、遗诏档、起居注册和内阁文书可以使死亡、即位与诏令生效的程序彼此校核；它们同样是由宫廷和官僚程序保留下来的记录，最擅长呈现何者被正式宣布，而不必然保存宗室、旗人和官员之间所有未公开的判断。康熙晚年继承程序缺乏公开而稳定的预期，足以解释信息为何被紧张地争夺，却不足以把每一种传闻都提升为证据。

正是在这种不确定性尚未完全散去时，朝廷要决定谁可以经手钱粮，谁须为亏空负责，何处的军需可优先动用。密折和朱批缩短了皇帝与个别督抚、将领之间的联络，但没有取消常规部院、地方衙门与驿路的作用。设置\eventref{EQ-1729-JUNJICHU-FOUNDING}{军机房（1729）}，是西北战争压力下为紧急军政事务另开较短的批答路径；它使皇帝更快看到呈报，也使催饷、调兵和问责可以并行。可是，快到北京的消息不等于快到前线的粮草。驼马、水源、季节和道路仍决定一项批示能否从案头转为行动。

\eventref{EQ-1735-QIANLONG-ACCESSION}{乾隆即位（1735）}时，继承因此不只是宫廷礼制的转换。新皇帝接手的是已经更密集的奏报和决策渠道，也是仍须按省、按案、按路线处理的财政和边务。雍正时期的档案连续，能说明官僚程序没有因改元而中断；不能由此断定前朝亏空已尽数厘清、地方社会已接受每项整顿，或新朝不再面对宗室、旗务与边地供给的压力。乾隆初政对于前朝人事、刑狱和边务的重新评估，正提醒人们：继位把统治的责任交接下来，却不会把既有的问题一并结案。

\subsection{省级中介与边地的不同时间}

中央能够看见的财政，须先由省级中介使之成为可呈报的财政。督抚、布政使、粮道、盐政官员、州县以及书吏和地方承办者，分别把田赋、漕粮、盐课、赈济和军需拆成不同的征解、留支与转运路线。对北京而言，一份奏销或题本可以把这些活动编成可审核的次序；对地方而言，同一项事务还牵连粮价、船只、仓储、民夫和信用。河工、漕运与盐政之所以反复进入高层讨论，恰恰因为富庶地区的税粮不是一笔可以不经转运而直接流入国库的收入。

这种行政结构在边地尤其显出不均。乾隆朝将西北远征转为长期经营时，\eventref{EQ-1760-ILI-GENERAL}{设伊犁将军（1760）}把驻防、屯田、安置和军台纳入一个军政枢纽。该项设置可以证明清廷企图让兵、地、粮与人事彼此衔接，不能证明伊犁河谷的牧地、北疆新城与南疆绿洲已经被同一种办法管理。军粮须沿道路转运，屯田须取得水、种子和劳力，驻军家庭、商人、原有居民与迁入者又面对不同的生计条件。将军衙门的发令和户籍册能够看见国家希望谁被安置在何处，却不能代替这些人是否留下、怎样取得土地和如何与邻近社群相处的记录。

\eventref{EQ-1765-USH-REBELLION}{乌什事件（1765）}使这一区别更难忽略。驻防城、军粮和官员任命并未使税役、劳役、地方权威与灌溉资源的关系自动稳定。清廷实录、军机档和方略能够追踪接报、调兵、镇压与善后的官方顺序；它们把冲突归入国家处理的问题时，也会压缩城中居民和地方中介所面对的具体压力。维吾尔或察合台文契约、地方文书与跨境商路材料若能与这些档案互读，才可能在个别地点追问水、地、贸易和征役怎样相互牵动。材料并不平均保存，故不能把一城的可见摩擦扩大为整个新疆的社会反应，也不能以「平定」一词取代后来治理实际展开的时间。

\subsection{财政的社会后果不止于国库}

高峰时期的财政秩序也通过身份进入家庭。八旗谱牒、出旗与改归民籍等安排，显示国家试图重排谁应由何种资源供养；名册却只记录分类和资格，不能直接显示一个家庭迁居后是否有地可耕、有工可做，或在新社区得到何种接纳。金川等战事中的夫役、牲畜、粮食与道路征用，则把军费化作村寨中劳力暂缺、市场受扰和债务风险的可能来源。攻克\eventref{EQ-1776-JINCHUAN-CONQUEST}{金川（1776）}结束了一轮军事行动，却没有替四川及邻省的运输者、纳户和战后迁徙者结清这些账目。

这里最容易被材料的形状误导。军机档、朱批奏折和方略保存得较多，因而朝廷何时发令、何人请饷、何处被报告为平定，往往比民夫如何离家、旗人家庭如何维生或屯田者如何分水更清楚。地方档案、账簿、契约、碑刻、寺院材料和不同语言的文书可以在局部补出后一类问题，但其保存偏向能写字、能订契约、能被衙门或寺院记录的人。档案的沉默不能证明妇女、佃户、临时迁徙者或未留下文字的社群没有承担代价。

因此，雍正、乾隆时期的「整顿」与「盛世」可以指向一种真实扩大的能力：朝廷较能把钱粮、报告、任命和军令接入同一条调度链。它们不应被写成对地方的全知全能。每一份账册都须问它记录的是定额还是实收，每一项命令都须问它是在北京发出、在省城转达，还是已在某个村寨、军台或绿洲实际施行。把这些层次分开，才能看见继承如何延续国家能力，也看见区域行政和社会生活如何不断限定那种能力的半径。
